% Complete reciprocal proofs for the separate finite-iteration packet.
\subsection{Finite-Rule Iterated Splicing}
\label{annex:p09-recombination-iterated-closures}

\begin{annexitem}{annex:p09.recombination-full-afl-iteration}{thm:p09-recombination-full-afl-iteration}{History nests and full-AFL closure}
\begin{proof}
If no rule is used, the closure is \(L\); this includes \(R=\varnothing\).
Thus assume some rule is used. The four bracket tokens have the meanings
declared in the reader. A balanced factor belongs to \(N\); in particular
an entire nest, a concatenation of nests and letters, and the empty word
are balanced. Every word in \(N\) has a unique decomposition into its
depth-zero letters and maximal nests: scan from left to right and match
each opening token to the first return to its starting depth, respecting
its rule label and its left/right type. This proves that erasing complete
maximal nests is well defined and that \(E(ab)=E(a)E(b)\) for \(a,b\in N\).

\paragraph{Lifting cuts and recording exactly the set closure.}
If \(E(x)=p\alpha_rq\), split the visible letter sequence immediately
before and after that occurrence of \(\alpha_r\). Attach each intervening
complete maximal nest to one neighboring piece. This gives
\(x=x_1vx_2\), with all three factors balanced and
\(E(x_1)=p,E(v)=\alpha_r,E(x_2)=q\).
This also works for an empty site by choosing a gap in the visible
sequence. No cut passes through an unmatched bracket. Do the same for
\(E(y)=p'\gamma_rq'\). The justified history operation then erases to
\(p\beta_rq'\), exactly the ordinary splice.

Conversely, the defining factorizations of every justified operation give
that same ordinary splice after applying \(E\).
Induction on a finite history construction therefore gives
\(E(J)\subseteq\sigma^*(L)\). For the opposite inclusion, use induction
on the stages \(\sigma^n(L)\). Choose histories for both input words;
the lifting just proved turns any ordinary splice into a justified history
whose erasure is its output. Inputs are retained, so different finite
stage indices can always be raised to their maximum. Thus
\(\sigma^*(L)\subseteq E(J)\).

Prune rules not used by any finite history construction. Every history
uses only retained rules, so this leaves \(J\) and \(E(J)\) unchanged.
Every retained rule has witnesses in \(J\) for both sites. No procedure
for deciding this pruning is needed for the existence argument.

\paragraph{Peeling a maximal left nest.}
If \(w=x\mathopen{(\!(}_ry]_rz\in J\), with the indicated nest maximal, then
\(xy\in J\) and \(\alpha_r\) is a prefix of \(E(y)\).
Prove this by induction on history length, equivalently on word length
along the construction. A last justified step has the form
\[
 w=p_1\mathopen{(\!(}_sp_2]_s\beta_s[_sq_1\mathclose{)\!)}_sq_2,
\]
with parents \(p_1p_2,q_1q_2\in J\), balanced factors, and
\(\alpha_s\) prefix of \(E(p_2)\), \(\gamma_s\) suffix of \(E(q_1)\).
There are exactly three locations for a maximal left nest:
the newly created \(\mathopen{(\!(}_sp_2]_s\), a maximal nest in \(p_1\),
or a maximal nest in \(q_2\). A nest inside the two new nests is not maximal.
In the new case peeling gives the first parent \(p_1p_2\).
In the \(p_1\) case, induction on that first parent gives \(xy\) directly.
In the \(q_2\) case, write \(q_2=x'\mathopen{(\!(}_ry]_rz'\).
Induction on the second parent gives \(q_1x'y\in J\) and the required
prefix condition. Reapply rule \(s\) to \(p_1p_2\) and \(q_1x'y\),
using the same site in \(q_1\); the result is \(xy\).
Both parents have smaller length, since a justified operation inserts
four bracket tokens and retains both parents inside its histories.

\paragraph{Exchanging the continuation of a maximal left nest.}
If \(x\mathopen{(\!(}_ry]_rz\) and \(x'\mathopen{(\!(}_ry']_rz'\) are in \(J\), with both nests
maximal, then \(x\mathopen{(\!(}_ry]_rz'\in J\).
Induct on the second history's construction. If its indicated nest is
new, \(r=s\) in its last-step expression above. Peeling the first history
gives \(xy\in J\), with \(\alpha_r\) prefix of \(E(y)\).
Splice it with the second history's second parent to recreate the retained
nest and that history's entire continuation \(z'\).
If the second indicated nest lies in \(q_2\), apply induction to its
second parent; the result already has the required final suffix \(z'\).
If it lies in \(p_1\), write its first parent as
\(x'\mathopen{(\!(}_ry']_rz''p_2\). Induction gives
\(x\mathopen{(\!(}_ry]_rz''p_2\in J\). Reapply the last rule \(s\) to this parent
and \(q_1q_2\). The unchanged suffix \(p_2\) still starts with its
required site after erasure, and the output is the required
\(x\mathopen{(\!(}_ry]_rz'\). These exhaust the possible maximal locations.

\paragraph{An explicit reversal isomorphism supplies the right cases.}
For each rule \(r\), define
\[
 r^\rho=\operatorname{rev}(u_4)\#
         \operatorname{rev}(u_3)\$
         \operatorname{rev}(u_2)\#
         \operatorname{rev}(u_1).
\]
Reverse word order, leave each letter unchanged, and send bracket tokens by
\[
 \mathopen{(\!(}_r\mapsto \mathclose{)\!)}_{r^\rho},\quad ]_r\mapsto [_{r^\rho},\quad
 [_r\mapsto ]_{r^\rho},\quad \mathclose{)\!)}_r\mapsto \mathopen{(\!(}_{r^\rho}.
\]
This involution sends left nests to right nests and satisfies
\(E(\rho(w))=\operatorname{rev}(E(w))\).
It sends a justified \(r\)-splice of \(x,y\) to a justified
\(r^\rho\)-splice of \(\rho(y),\rho(x)\): substituting the four
factor reversals in the defining history word verifies every position.
The new sites and recombinant are respectively
\(\operatorname{rev}(\gamma_r),\operatorname{rev}(\alpha_r)\),
and \(\operatorname{rev}(\beta_r)\).
Hence it bijects \(J\) with the history language for the reversed initial
language and transformed rules.

Applying the preceding proved statements through this involution gives,
without an omitted dual argument:
if \(x[_ry\mathclose{)\!)}_rz\in J\) is maximal, then \(yz\in J\) and
\(\gamma_r\) is a suffix of \(E(y)\);
and two maximal right nests of the same label allow replacement of the
prefix of one by the prefix of the other while keeping its nest and suffix.
This is a structural isomorphism, not an assumption that \(FA\) is closed
under reversal.

\paragraph{Promoting a bracket adjacency.}
If \(bqb'\), with \(q\in V^*\) and \(b,b'\) bracket tokens, occurs in
a history, choose a shortest history containing this occurrence.
At least one endpoint is an endpoint of a maximal nest.
Otherwise both endpoints lie inside the same maximal nest: since the
intervening word contains no brackets, they cannot belong to distinct
maximal nests without an intervening endpoint. Peel that containing nest,
using the left or right peeling result. The occurrence is preserved inside
the peeled content and the history becomes shorter, a contradiction.
In particular an occurrence \(\mathopen{(\!(}_sq\mathopen{(\!(}_t\), or \(\mathopen{(\!(}_sq[_r\), can be witnessed
with its \(\mathopen{(\!(}_s\)-nest maximal: the later opening token is still inside
that unclosed nest, so it cannot itself be a maximal opening.

\paragraph{Regular prefix and suffix descriptors.}
For fixed \(r,s,t\), each short occurrence set in \(F_r\) is a subset of
the finitely many words of length below \(|\alpha_s|\).
The terminal remainder \(\alpha_sV^*\) is regular.
Use a finite graph with vertices \(R\cup\{\$\}\) and an edge \(s\to t\)
labeled by the regular language \(F_r(s,t)\).
Replace each edge by a finite automaton for its label, using disjoint
internal states and epsilon edges to its endpoints. Accept paths from
the chosen \(s\) to \(\$\). This finite automaton accepts exactly
\(S_r(s)\), including zero-length edge labels and arbitrary finite paths.
Set the dummy entry to \(\{\lambda\}\).

Perform the same construction for the reversed history system.
Define \(P_r(s)\) to be the reversal of
\(S'_{r^\rho}(s^\rho)\), with \(\$^\rho=\$\).
Reversing an automaton's arrows and interchanging its initial/final
conditions, with a fresh start if needed, proves these prefix sets regular.
This defines the full dual operation explicitly.
The family hypothesis is only used for quotients and substitution below.

\paragraph{Tracing visible blocks: \(E(J)\subseteq K\).}
A justified operation retains a balanced prefix of its first parent,
inserts a new left nest, the plain word \(\beta_r\), a new right nest,
and a balanced suffix of its second parent.
Consequently no two old plain blocks merge: there is a bracket between
the retained pieces and the inserted block. Every depth-zero plain block
is therefore a contiguous interval of a unique birth block, either a
word of \(L\) or a newly inserted \(\beta_r\). Track a gap as well as
letters when the interval is empty; the same description then applies
to empty visible blocks. Give a birth word of \(L\) index \(\$\),
and a birth word \(\beta_r\) index \(r\).

Here is the explicit right-side tracing for such an interval \(y\).
If its right boundary is the original \([_r\), or the end of its
dummy-index birth word, take \(t=\$\) and suffix \(\lambda\).
Otherwise that boundary is \(\mathopen{(\!(}_t\), from a later truncation.
Peel this maximal left nest. The discarded side starts, after erasure,
with \(\alpha_t\). Follow the truncations of this same birth interval
backwards, along the discarded side. Between consecutive opening tokens
\(\mathopen{(\!(}_{s_i},O_{s_{i+1}}\) the plain gap \(q_i\) has one of two roles.
If it already starts with the whole \(\alpha_{s_i}\), stop: the
entire remaining suffix of the birth word is in
\(F_r(s_i,\$)=\alpha_{s_i}V^*\).
If it does not, peeling's prefix condition implies
\(|q_i|<|\alpha_{s_i}|\): a longer raw initial segment would already
contain the whole required site. The encountered adjacency is
\(\mathopen{(\!(}_{s_i}q_iO_{s_{i+1}}\), in a history or its peeled ancestor,
so \(q_i\in F_r(s_i,s_{i+1})\).
At a non-dummy birth boundary \([_r\), the last short gap instead
witnesses \(\mathopen{(\!(}_{s_i}q_i[_r\), giving the short terminal alternative.
At a dummy birth end there is no further bracket to supply a missing site,
so the ordinary terminal alternative must apply.

Each backward step removes a truncating history nest and reaches an earlier
finite construction. Thus this process terminates. The starting index is
\(s_1=t\), adjacent indices match, and the final target is \(\$\);
the suffix \(q=q_1\cdots q_n\) belongs to \(S_r(t)\).
This procedure traces to the original birth block, not to a guessed
arbitrary completion. The only right-side markers encountered before
that birth boundary are opening markers of successive right truncations;
the original boundary is \([_r\). This follows inductively from the
balanced-prefix/suffix form of the history operation above.
Reversal gives the left-side tracing, with discarded prefix
\(p\in P_r(s)\). The two traces refer to the same birth interval, so
their restored birth word is \(pyq\in L_r\).
Hence \(y\in L_r(s,t)\).

For clarity the four internal configurations give precisely these labels.
Between \(\mathopen{(\!(}_sx]_s\) and \([_tz\mathclose{)\!)}_t\), the unchanged birth frame has
\(s=t\) and the visible word is \(\beta_s\), hence \(L_s(\$,\$)\).
Between a left nest and a left nest, the left boundary retains the
birth frame's \(]_r\) and the right trace gives \(L_r(\$,t)\).
Between two right nests the reversed argument gives \(L_r(s,\$)\).
Between a right and a left nest both traces are present, giving
\(L_r(s,t)\). In this last case the original birth boundary inside the
left-side history identifies \(r\): it is dummy for a word born in \(L\),
or the label of that birth's square closing token. The right-side trace
identifies the same birth word. This also proves consistency of the two
traces rather than assigning unrelated quotient witnesses.

The first visible block has dummy birth and no discarded prefix;
the last has dummy birth and no discarded suffix.
At a separating left nest, its label is simultaneously the preceding
block's \(t_i\) and the following birth's \(r_{i+1}\), whose left
quotient index is dummy. At a separating right nest, the preceding
right quotient is dummy and its birth index is the following \(s_{i+1}\).
These are exactly the two adjacency alternatives defining \(K\).
Deleting maximal nests concatenates these compatible blocks, proving
\(E(J)\subseteq K\).

\paragraph{A uniform right-trimming construction.}
Suppose \(q\in S_r(t)\), \(t\in R\), and
\(XbqZ\in J\), where \(b,q\) are plain words and \(X,Z\) are balanced.
Require \(Z=\lambda\) if \(r=\$\), or \(Z\) to start with \([_r\)
if \(r\in R\). Then there exists balanced \(Z'\) starting with \(\mathopen{(\!(}_t\)
such that \(XbZ'\in J\). The entire opposite context \(Xb\) is preserved.

Choose a path factorization of \(q\) and induct on its number of edges.
For one edge, \(q\in F_r(t,\$)\).
If \(\alpha_t\) is a prefix of \(q\), choose an available second history
\(a_0v_0b_0\) with \(E(v_0)=\gamma_t\).
A justified splice with first prefix \(Xb\) wraps the whole remainder
\(qZ\), giving
\[
 Xb\mathopen{(\!(}_t(qZ)]_t\beta_t[_ta_0v_0\mathclose{)\!)}_tb_0.
\]
Its new balanced suffix starts with \(\mathopen{(\!(}_t\), as required.

In the short terminal case \(r\in R\), write
\(Z=[_rA\mathclose{)\!)}_rB\).
Promote a witness for \(\mathopen{(\!(}_tq[_r\) to get
\[
 X'\mathopen{(\!(}_tq[_rU\mathclose{)\!)}_rV]_tW\in J
\]
with its \(\mathopen{(\!(}_t\)-nest maximal. Peeling gives
\(X'q[_rU\mathclose{)\!)}_rV\in J\), with \(\alpha_t\) prefix of
\(E(q[_rU\mathclose{)\!)}_rV)\).
The right-nest exchange proved above combines this history with
\(Xbq[_rA\mathclose{)\!)}_rB\), producing
\(Xbq[_rU\mathclose{)\!)}_rV\in J\).
Splice it with the available \(\gamma_t\)-history, taking prefix \(Xb\);
the resulting suffix is
\(\mathopen{(\!(}_t(q[_rU\mathclose{)\!)}_rV)]_t\beta_t[_ta_0v_0\mathclose{)\!)}_tb_0\).
Every required site is justified by the displayed prefix condition.

For more than one edge, write \(q=q_1q'\), with
\(q_1\in F_r(t,u)\) and \(q'\in S_r(u)\), \(u\in R\).
Induction supplies
\(Xbq_1\mathopen{(\!(}_uA]_uB\in J\).
Promote a witness for \(\mathopen{(\!(}_tq_1\mathopen{(\!(}_u\); it has the form
\(X'\mathopen{(\!(}_tq_1\mathopen{(\!(}_uU]_uV]_tW\).
Peeling gives \(X'q_1\mathopen{(\!(}_uU]_uV\in J\), with \(\alpha_t\)
prefix of \(E(q_1\mathopen{(\!(}_uU]_uV)\).
Left-nest exchange preserves the first history's prefix and \(\mathopen{(\!(}_u\)-nest
while replacing its continuation by \(V\), yielding
\(Xbq_1\mathopen{(\!(}_uA]_uV\in J\).
Its discarded side has the same erasure
\(E(q_1\mathopen{(\!(}_uA]_uV)=E(q_1\mathopen{(\!(}_uU]_uV)\), since both \(\mathopen{(\!(}_u\)-nests
erase entirely. A justified \(t\)-splice therefore wraps this side,
again preserving \(Xb\) and producing a suffix starting with \(\mathopen{(\!(}_t\).
This completes the induction, including empty short gaps.

By the explicit reversal isomorphism, the exact left-trimming assertion
is also proved: if \(p\in P_r(s)\) and \(XpbZ\in J\), with
\(X=\lambda\) for dummy \(r\) or ending with \(]_r\) otherwise,
then a history \(X'bZ\in J\) exists with \(X'\) ending with \(\mathclose{)\!)}_s\).
It preserves the entire opposite context \(bZ\).

\paragraph{Every quotient piece has the required history boundaries.}
Let \(y\in L_r(s,t)\) and choose \(p,q\) with \(pyq\in L_r\).
For \(r=\$\), the word \(pyq\) itself is a history, with empty contexts.
For \(r\in R\), availability of both sites permits a justified \(r\)-splice;
it yields a history \(X(pyq)Z\), where \(X\) ends with \(]_r\)
and \(Z\) starts with \([_r\), since \(pyq=\beta_r\).
If \(t\ne\$\), right trimming gives \(XpyZ_t\), with \(Z_t\)
starting with \(\mathopen{(\!(}_t\); otherwise \(q=\lambda\) and retain \(Z\).
If \(s\ne\$\), apply left trimming to this very history.
Its opposite context is now \(yZ_t\), which the lemma preserves.
This gives \(X_syZ_t\) with \(X_s\) ending with \(\mathclose{)\!)}_s\).
If \(s=\$\), \(p=\lambda\) and no left trimming is necessary.
Thus the two-sided case is proved with actual commuting context-preserving
operations, not an assumed double induction.

The resulting prefix is empty when \(s=r=\$\), ends with \(]_r\)
when \(s=\$,r\ne\$\), and ends with \(\mathclose{)\!)}_s\) when \(s\ne\$\).
The suffix is empty when \(t=r=\$\), starts with \([_r\)
when \(t=\$,r\ne\$\), and starts with \(\mathopen{(\!(}_t\) when \(t\ne\$\).
These are exactly the source's six boundary conditions.

\paragraph{Assembling compatible pieces: \(K\subseteq E(J)\).}
Choose a defining concatenation \(y_1\cdots y_n\) in \(K\).
The boundary result gives a history \(y_1Z_1\), since
\(r_1=s_1=\$\). If \(n=1\), also \(r_1=t_1=\$\), so it is just
\(y_1\in J\). Otherwise its suffix has the required first opening token.

Inductively suppose \(B_kZ_k\in J\), both factors balanced, with
\(E(B_k)=y_1\cdots y_k\), and with suffix starting with \([_{r_k}\)
if \(t_k=\$,r_k\ne\$\), or \(\mathopen{(\!(}_{t_k}\) if \(t_k\ne\$\).
Choose \(Xy_{k+1}Z\in J\) from the boundary result.
In the left-nest adjacency case \(t_k=r_{k+1}=r\ne\$\) and
\(s_{k+1}=\$\), write
\(Z_k=\mathopen{(\!(}_rA]_rB\) and \(X=X'\mathopen{(\!(}_rU]_r\).
Left-nest exchange produces
\(B_k\mathopen{(\!(}_rA]_ry_{k+1}Z\in J\).
Take \(B_{k+1}=B_k\mathopen{(\!(}_rA]_ry_{k+1}\).
In the right-nest adjacency case \(t_k=\$,r_k=s_{k+1}=s\ne\$\),
write \(Z_k=[_sA\mathclose{)\!)}_sB\) and \(X=X'[_sU\mathclose{)\!)}_s\).
Right-nest exchange produces
\(B_k[_sU\mathclose{)\!)}_sy_{k+1}Z\in J\).
Take \(B_{k+1}=B_k[_sU\mathclose{)\!)}_sy_{k+1}\).
In both cases the inserted nest erases, so
\(E(B_{k+1})=E(B_k)y_{k+1}\), and the new suffix \(Z\)
has the boundary required for the next step.
At the last step \(r_n=t_n=\$\) and \(Z=\lambda\).
Thus \(B_n\in J\) erases to the entire chosen concatenation.
This proves \(K\subseteq E(J)\).

\paragraph{Membership in the family.}
All \(P_r,S_r\) are regular. Each \(L_r\) is either \(L\in FA\)
or a singleton, so successive regular quotients give
\(L_r(s,t)\in FA\).
There are finitely many triples \((r,s,t)\). A finite automaton over
these triple labels enforces the initial, terminal, and adjacency conditions:
its states record the last triple, with transitions precisely for the
two allowed alternatives. Substitute \(L_r(s,t)\) for each label.
The family's regular-substitution hypothesis gives \(K\in FA\).
The two established inclusions and \(E(J)=\sigma^*(L)\) finish the proof.
No arbitrary increasing-union closure, reversal closure of \(FA\),
or consuming-multiset assumption was used.
\end{proof}
\end{annexitem}

\begin{annexitem}{annex:p09.recombination-regular-iteration}{thm:p09-recombination-regular-iteration}{Regular-language specialization}
\begin{proof}
Regular languages satisfy the three stated properties.
For a right quotient \(L/A\) with both languages regular, run the automaton
for \(L\) on the visible input and accept at exactly those states from which
some word accepted by an automaton for \(A\) reaches a final \(L\)-state.
These states are determined by finite product-graph reachability.
For a left quotient \(A\backslash L\), use as initial states the
\(L\)-states reached by words in \(A\), again a finite product computation,
and use the original \(L\)-final states. These constructions handle the
empty word.

For regular-control substitution, replace every transition labeled \(a\)
in the control automaton by a disjoint finite automaton for its substituted
regular language, joining entry and exit by epsilon edges.
The resulting finite automaton accepts exactly all allowed concatenations.
Thus the full-AFL theorem applies to regular \(L\) and finite \(R\).
It gives \(H(REG,FIN)\subseteq REG\).
For the reverse inclusion choose the empty rule set, which leaves each
regular input unchanged. Hence equality holds.
\end{proof}
\end{annexitem}

\begin{annexitem}{annex:p09.recombination-radius-hierarchy}{thm:p09-recombination-radius-hierarchy}{Exact radius witness and strict endpoints}
\begin{proof}
Fix \(k\geq1\), let \(A=a^{2k}b^{2k}\), \(B=b^{2k}a^{2k}\),
and use the source's seed \(Aa^{2k+2}B\) and rule
\(a^{k+1}\#a^k\$a^{k+1}\#a^k\).
Both recognition sites have length \(2k+1\); the exterior \(a\)-blocks
have length only \(2k\), so sites occur exclusively in the central block.
For inputs \(Aa^nB,Aa^mB\), cuts after central positions \(i,j\) satisfy
\(k+1\leq i\leq n-k\), \(k+1\leq j\leq m-k\).
The output is \(Aa^{i+m-j}B\), whose central length is at least \(2k+1\).
Thus every generated word lies in \(L_k\).
Two copies of the seed, with \(i=k+1,j=k+2\), give length \(2k+1\).
Given central length \(n\geq2k+1\), splice that word with the seed
using \(i=n-k,j=k+1\); the new central length is \(n+1\).
Induction gives every required length, proving exact generation.
The maximal handle length is \(k+1\).

Suppose a radius-at-most-\(k\) scheme and finite seed generated \(L_k\).
Then \(L_k\) is closed under its rules. Any applicable rule has left site
\(u_1u_2\) of length at most \(2k\). On a word of \(L_k\), such a site
cannot cross two changes of letter: that would require at least
\(2k+2\) letters, since every intervening whole run has length at least \(2k\).
It therefore has one of the following forms. Keep any applicable second
input and its retained suffix fixed, and move the first site's occurrence:
\begin{enumerate}
\item For a site in \(a^*\), including the empty site, choose matching
      occurrences at the starts of the central and final \(a\)-blocks.
      The retained prefixes contain respectively \(2k\) and \(4k\)
      occurrences of \(b\), regardless of the cut inside the site.
\item For a nonempty site in \(b^*\), use the same offset in the first
      and second \(b\)-blocks. The retained-prefix \(b\)-counts differ by \(2k\).
\item For a site in \(a^+b^+\), use its matching offsets across the
      first \(a/b\) boundary and the central \(a/b\) boundary.
      Its \(a\)-part and \(b\)-part each fit; the retained-prefix
      \(b\)-counts again differ by \(2k\).
\item For a site in \(b^+a^+\), use the first \(b/a\) boundary and
      the second \(b/a\) boundary. The same fixed cut gives
      retained-prefix \(b\)-counts differing by \(2k\).
\end{enumerate}
These exhaust all sites. In each case both chosen retained prefixes are
nonempty: in the mixed \(a^+b^+\) case the \(a\)-part has length at most
\(2k-1\), so its first occurrence starts after at least one leading \(a\).
Every word in \(L_k\) has exactly \(4k\) letters \(b\).
Joining the fixed suffix to prefixes with different \(b\)-counts therefore
gives at least one nonempty output outside \(L_k\), contradicting closure.
Thus no rule can be applicable on \(L_k\) in an exact generating scheme.
Its iterates would leave its finite seed unchanged, whereas \(L_k\) is infinite.
This proves \(L_k\notin H(FIN,k)\), including the source's convention
of ignoring empty-word differences.

Radius monotonicity gives every displayed inclusion between adjacent
\(H(FIN,k)\)'s; the witnesses just proved make them strict.
Every finite language is generated unchanged by an empty-rule scheme.
The radius-one rule \(a\#\$\#a\) from seed \(\{a\}\) generates exactly
\(a^+\): every output is a nonempty unary word, and retaining the whole
first input and the one-letter second input increases its length by one.
This proves the first strict inclusion.
Every finite scheme has some finite radius, so the finite-rule family
contains all the radius families; each fixed radius family is proper in
it because of its next-radius witness. The regular-iteration theorem
puts the whole finite-rule family inside \(REG\).

For the final separation let
\(W_1=(dc)^*a^+(dc)^*\),
\(W_2=c(dc)^*a^+(dc)^*d\), and \(W=W_1\cup W_2\).
This is an infinite regular language, and \(W_1,W_2\) are disjoint:
the first type begins with \(d\) or \(a\) and ends with \(c\) or \(a\);
the second begins with \(c\) and ends with \(d\).
For every \(z\in W\), \(cz\notin W\):
on \(W_1\) its ending excludes \(W_2\) and its beginning excludes \(W_1\);
on \(W_2\) its first two letters are \(cc\), excluded by both types.
Likewise \(dz\notin W\): on \(W_1\) it begins with \(da\) or \(dd\),
and on \(W_2\) its ending is \(d\) although its beginning is not \(c\).

If an oriented splice of \(x,y\in W\) yielded \(z\in W\), and \(x\in W_1\),
replace the first input by \(cxd\in W_2\), retaining the same recognition
site and second input. Its output is \(cz\notin W\).
If \(x\in W_2\), replace it by \(dxc\in W_1\), giving \(dz\notin W\).
Thus a scheme closed on \(W\) can have no applicable rule on \(W\).
It cannot generate this infinite language from a finite seed.
The argument does not assume anything about the complexity of the rule set;
in particular \(W\notin H(FIN,FIN)\).
This proves the final strict inclusion with the source's actual wrapper
mechanism, and completes the hierarchy.
\end{proof}
\end{annexitem}

\begin{annexitem}{annex:p09.recombination-nest-domain}{thm:p09-recombination-nest-domain}{Context-free and nonregular history domain}
\begin{proof}
Use one grammar nonterminal T, with the finite productions, stated inline:
\(T\to\lambda\), \(T\to aT\) for each \(a\in V\),
\(T\to \mathopen{(\!(}_rT]_rT\), and \(T\to [_rT\mathclose{)\!)}_rT\) for each \(r\in R\).
Induction on derivations shows that every generated word is well nested.
Conversely, a nonempty well-nested word begins either with a letter,
or with a complete left or right nest followed by another well-nested word.
Its unique first matching endpoint gives exactly one of these decomposition
forms. Induction on word length supplies derivations for the shorter
contents and suffix. Thus the grammar generates exactly N, proving
context-freeness, including empty contents and empty suffixes.

If R is nonempty, choose a label r. Intersect N with the regular language
\(\mathopen{(\!(}_r^*]_r^*\). Only left nests of this label can occur, and all their
openers precede all closers, so the intersection is exactly
\(\{\mathopen{(\!(}_r^n]_r^n:n\geq0\}\).
Were N regular, its finite-automaton product with an automaton for this
intersection condition would be a finite automaton for that language.
For an m-state deterministic automaton, consider the accepted word
\(\mathopen{(\!(}_r^{m+1}]_r^{m+1}\). A repeated state among the first m+1 opener
prefixes gives a nonempty opener-only cycle. Repeating it once adds
openers but no closers and leaves acceptance unchanged, contradicting
the required equality of counts. Hence N is not regular.
If R is empty, there are no bracket tokens and N equals \(V^*\).
\end{proof}
\end{annexitem}
